\section{Discussion}

\paragraph{Why factored residuals help.}
The dense residual baseline has enough capacity to fit cross terms in source data, but the same flexibility can mix city-specific demand, speed, and headway effects.
\ours restricts each output to a smaller mechanism-specific parent set.
This constraint appears helpful in the strict cross-city setting: it reduces total MSE on all four targets and makes source weighting more effective.
The formal risk decomposition in the method section explains this as a bias-control tradeoff: mechanism-wise models reduce cross-mechanism mixing when transfer shifts are localized, while dense residuals may still be competitive when the selected sparse features are the main source of gain.
These experiments support mechanism-factored transfer as a useful inductive bias; they do not establish formal causal identification or invariance under intervention.

\paragraph{What the source weighting and source selection do and do not show.}
Similarity weighting improves the mean strict leave-one-out ratio from $0.759$ to $0.578$ and wins over unweighted \ours on three of four targets.
The Singapore target is the exception, where weighting is slightly worse than unweighted \ours.
The full-module feature analysis makes the boundary sharper: a source-only feature-set selector chooses simulator-output summaries for most splits, observation summaries for the duplicated Singapore target, and never selects local-only or uniform weighting.
Thus the useful signal in this diagnostic is not the weak static local graph itself, but the target's unlabeled transition-summary statistics.
An expanded NYC/SEPTA/MBTA GTFS-RT occupancy snapshot check gives the same warning: static local features have cross-city mean ratio $1.011$ against no-local features and the selector falls back to no-local, while within-city temporal prediction is essentially tied with mean ratio $0.997$.
The guarded full-module ensemble makes this boundary more explicit.
It helps when the unlabeled target summary identifies a dominant nearby source, as in MBTA and Austin, but the guard falls back to rigid \ours when source similarity is diffuse, as in Halifax.
The same-information dense \Htwo source ensemble also improves substantially, so the target-summary signal is genuinely useful; the remaining gap favors mechanism-factored residuals over dense residual mixing.
This is not target-label calibration: the target labels remain unseen, but the method is not a pure source-only setting because it uses unlabeled target summaries for source selection.

\paragraph{Zero-shot is a lower bound, not the intended deployment endpoint.}
The single-source transfer matrix is useful precisely because it is uneven.
It shows that a mechanism model trained on one city can transfer in some directions, but it also exposes source-target pairs where one-city history is not enough.
The stronger and more realistic claim is a target-information ladder: multi-source target-label-free transfer gives the main zero-target-label result, and few-shot target offline adaptation uses a small amount of passive target history to tune mechanism trust without claiming online deployment.
This framing matches practical transit deployment better than a pure zero-shot claim, because agencies normally have at least some archived schedules, counts, AVL traces, or APC samples before changing a control policy.
The Austin APC time-budget sweep makes this concrete: one hour is too little to trust, six hours is enough for the residual gate to reach parity or better in both slices, and one day to one week gives a clearer target-adaptation benefit.
The target-route baseline comparison also changes the preferred few-shot implementation.
Full mechanism refitting is not the safest way to use a small target budget; a low-dimensional target context-bias adapter is consistently stronger on average.
This is consistent with the causal-mechanism view: passive target history should first update mechanism state or trust, not overwrite all transferred mechanism parameters.

\paragraph{Policy evidence remains secondary.}
The one-step policy proxy and sampled BusSimEnv rollout are useful checks, but they do not replace a closed-loop operational benchmark.
The expanded sampled rollout is especially important because it weakens the earlier parity story: the raw aggregate favors \Htwo, while city-normalized headway and bunching metrics favor or nearly tie weighted \ours.
That correction is useful rather than damaging; it keeps the policy claim bounded and prevents the paper from overstating a dynamics-transfer result as operational dominance.
The full-line libsumo diagnostic strengthens the control story only after adding a safe target-summary prior.
Unconstrained learned residual MPC improves over dense learned MPC but remains far weaker than classical holding rules in severe bunching states.
By contrast, a target-summary gated Daganzo prior, which uses static route and schedule summaries to choose the holding gain, slightly improves over the best fixed Daganzo gain in the four-city diagnostic.
This suggests that the next controller should not be a free-form residual action selector.
It should instantiate a target mechanism state from route, schedule, demand, and AVL/APC summaries, use that state to gate a physically meaningful holding prior, and allow learned residual actions only behind a trust gate.
The strongest claim should therefore remain dynamics transfer plus bounded policy improvement, not field-ready policy dominance.

\paragraph{What the RESCO stress test changes.}
The RESCO result gives a cleaner counterfactual-control test than passive bus logs because each candidate signal phase can be rolled out in SUMO.
It also imposes a harder baseline: MaxPressure is close to optimal in several regular network regimes.
The pressure-guarded CFCMT controller should therefore be interpreted as a safe transfer design, not as a universal replacement for pressure control.
On the extended RESCO subset it reaches MaxPressure-level mean queue, is close to phase pressure with a non-significant mean advantage, and significantly improves over dense residual transfer, simulator-only MPC, spillback pressure, and fixed programs.
The 3600-second check is even more conservative: phase pressure becomes clearly best, while guarded CFCMT remains closer to MaxPressure and still better than dense residual transfer and fixed programs.
The non-significant short-horizon gaps against MaxPressure and phase pressure, together with the long-horizon phase-pressure win, are scientifically useful: they show that the method must preserve physically meaningful priors when a classical rule is already near the ceiling.
This is also the most defensible direction for future policy work in bus holding, where a gated Daganzo or threshold prior is currently stronger than unconstrained learned residual MPC.

\paragraph{Why the RL leaderboard tables are separated.}
The official RESCO and LibSignal full-budget runs are included to address a common traffic-control review question: whether the proposed controller was compared against standard RL baselines rather than only hand-coded rules.
They also expose a practical point.
On these benchmarks, full-budget RL training is expensive and implementation-specific; for example, LibSignal \texttt{FRAP} on \texttt{sumo4x4} is orders of magnitude slower than small-grid runs even after cProfile-driven simulator and parallelization fixes.
More importantly, these libraries do not use the same observation representation, reward, control interval, or metric aggregation as the \ours phase-MPC stress test.
Putting all rows into one ranked table would therefore create a false sense of protocol equivalence.
The correct use is as a native-protocol leaderboard cross-check in the appendix, while the main causal-transfer comparison remains the same-protocol phase-MPC table.

\paragraph{Static-derived validation is not operational validation.}
The central limitation is target construction.
The experiments use full-route open static data and deterministic transition rules to create paired simulator/target transitions.
This is valuable for reproducibility and reviewer scrutiny, but it is not the same as evaluating counterfactual holding actions on AVL/APC trajectories or a calibrated live SUMO rollout.
The Austin raw APC/AVL slice partially addresses the trajectory-ground-truth concern for passive action-$0$ dynamics, but it still cannot identify how unobserved holding actions would have changed passenger and headway outcomes.
True counterfactual holding labels would require an intervention design: randomized or rule-threshold hold/no-hold variation in the field, a natural experiment that changes holding behavior while leaving comparable states, or a simulator calibrated and validated well enough to serve as the counterfactual environment.
The paper should therefore avoid claims such as ``CFCMT is better on real bus systems''.
A precise claim is: under open static-data-derived cross-city validation with no or few target labels, and under passive real APC validation with few-shot adaptation, mechanism-factored residual transfer with unlabeled target-summary source selection is more robust than dense residual correction.

\paragraph{How to strengthen the next version.}
The most valuable addition would be counterfactual or interventional validation on one city with vehicle-level AVL/APC trajectories, even if only for a subset of routes.
The second priority is a calibrated or independently validated closed-loop simulator, so the target-summary-gated prior and any learned residual override can be tested against real or validated counterfactual outcomes.
The third priority is to add more cities, because city-level generalization is ultimately limited by the number of held-out cities.
